\documentclass[11pt,a4paper,titlepage]{article}

\title {Alagang MMC: A System to Address The Process Hurdles in the Makati Medical Center's Patient Care Coordination Leveraging Human-Centered Design Through The CeHRes Roadmap}
\author {Augustus Clark Raphael P. Rodriguez \\ Department of Physical Sciences and Mathematics \\ University of the Philippines Manila \\ \\
	Adviser: Marbert John C. Marasigan, M.Sc.}
\date{}

% \setlength{\parindent}{2em}   % first-line indent size


\usepackage[top=1in, bottom=1in, left=1.5in, right=1in]{geometry}
% CUSTOM PACKAGES, USE IN FINAL DOCUMENT
\usepackage{graphicx}
\usepackage{setspace}
\usepackage{xcolor}
\usepackage{rotating}
\usepackage{pdflscape}

% \usepackage[hidelinks]{hyperref}
\doublespacing

% allows for separate page section by changing the behavior of \section
\let\stdsection\section
\renewcommand\section{\newpage\stdsection}

\begin{document}
	\maketitle
	\tableofcontents
	%  KIOSK IS JUST A DUMB TERMINAL WITH STATIC DATA KIOSK IS JUST A DUMB TERMINAL WITH STATIC DATA KIOSK IS JUST A DUMB TERMINAL WITH STATIC DATA KIOSK IS JUST A DUMB TERMINAL WITH STATIC DATA KIOSK IS JUST A DUMB TERMINAL WITH STATIC DATA KIOSK IS JUST A DUMB TERMINAL WITH STATIC DATA KIOSK IS JUST A DUMB TERMINAL WITH STATIC DATA KIOSK IS JUST A DUMB TERMINAL WITH STATIC DATA KIOSK IS JUST A DUMB TERMINAL WITH STATIC DATA KIOSK IS JUST A DUMB TERMINAL WITH STATIC DATA 
	\section{Introduction}
	\subsection{Background of the Study}
	The Makati Medical Center (MMC) is located at the heart of the Philippines' business district and offers numerous healthcare services paired with modern, state of the art facilities. Accredited by the Department of Health, United Nations Children's Fund, World Health Organization, the Joint Commission International, alongside a number of organizations, the institute and its highly skilled and board-certified staff work under a commitment to improving patient health outcomes in the nation \cite{makatimedicalcenter_about, makatimedicalcenter_2025_makatimed}. The hospital is among the leading healthcare institutes in the country, with local articles discussing its diverse capabilities as an institute. With a long-standing reputation of 54 years in service, the hospital continues to provide quality healthcare to its patients while minimizing waiting times as much as possible \cite{kritikakritika_2025_top, metroguide_2021_the, flor_2023_top, makatimedicalcenter_history}. 
	%\textcolor{red}{the leading healthcare institute}
	
	% Set the scene and the background of your SP topic.
	Digital Health in accordance with the United States Food Drug Administration \cite{foodanddrugadministration_2020_what} is the group of broad categories encompassing technologically adjacent health domains. The value of a digital health system lies in its potential to drive new policies, improve care quality, better educate consumers, and most especially decrease any inflationary costs. As the Philippines continues to adopt more technologies, it is important to consider how these technologies can be adopted into the country's health infrastructure. Sucaldito et al. \cite{sucaldito_2025_challenges} conducted a review of health information systems (HIS) currently implemented and found low scores in governance indicative of deficiencies in the HIS implementation strategy citing that differences in infrastructure and funding need to be addressed or worked around to improve these implementations. Reinforcing the need for continued developments in the field in order to address gaps in critical areas. 
	
	Guzman \cite{guzman_2025_first} writes that the Philippines in recent times has been able to adopt systems that connect virtual appointments, telehealth consultations, hospital diagnostics, and electronic medical records (EMR) into a single system for the convenience of not only the hospital, but also the patients . Such developments demonstrate the feasibility of producing digital health systems in the country despite the challenges that may arise. The rapid growth of technology has a proportional impact on the development of digital health systems. Livieri et al. \cite{livieri_2025_the} found that among the top challenges that arise in development stem from generational gaps followed by ethical concerns and mistrust in the technology. This suggests Digital Health is a continuously evolving domain that evidently does not have a ``one-size-fits-all" approach and a well-informed implementation strategy is essential in addressing technological gaps. 
	
	Reflective of the national trend towards integrated digital healthcare, MMC has established comprehensive digital health systems within its operations and among its clinicians.  In a personal interview with Dr. Andrei Rodriguez \cite{Rodriguez_A}, they discuss the history of MMC and how it was the first hospital in the Philippines to successfully transition to a digitized environment via its Electronic Medical Record (EMR) project in conjunction with the hospital's integrated Hospital Information System (iHIMS). Further articles discuss that as early as 2013, the hospital was the first in the Philippines to partner with MIMS to integrate MIMS gateway into iHIMS, enabling clinicians to reduce medication errors and ensure correct prescriptions through accurate medication safety alerts \cite{mimspteltd_2016_mims,makatimedicalcenter_2016_makati}.
	
	Human-Centered Design (HCD) is a rigorous approach that primarily puts stakeholders' needs first in the design and development processes of a user-facing system \cite{leary_2022_an}. The philosophy has seen recent successes in the healthcare domain, for instance, the work done by Landeiro et al. \cite{landeiro_2025_humancentered} with regards to enhancing a mother's experience during the maternity journey. The authors discuss how the approach enabled tailored solutions for personalized care and improving the non-clinical aspects of the maternity journey, leading to enhanced care experience and humanization, supported emotional well-being and mental health, increased empowerment and communication, and addressed a high-risk and vulnerable population. 

	Maximizing the continued adoption of newer technologies in the Philippines in order to progressively digitize core workflows may enhance overall workplace efficiency over manual record management and book-keeping. Libadia et al. \cite{libadia_2025_development} found that a reliable system that integrates seamlessly into the current workflow and offers relevant features boosted workplace compliance and productivity . Their approach put the user first and addressed any existing logistical burdens and other such burdens caused by outdated measures. The development of a support system geared towards streamlining the patient-doctor relationship may provide ample opportunities to help the hospital continue to keep up with the times and diversify patient reach. 
	
	This study aims to design and develop a digital health intervention for the MMC grounded on HCD, ensuring that the stakeholders needs are met in every phase of development. The proposed system aims to address burdens the institute faces with regards to its internal systems and processes.
	
	
	\subsection{Statement of the Problem}
	According to Dr. Andrei Rodriguez, Chair of the Peer Review Committees at MMC and Rheumatologist, and Dr. Saturnino Javier, current Medical Director and Cardiologist \cite{Rodriguez_A, Javier_S}, the MMC's patient care coordination measures are burdened with process hurdles due to the outdated nature of existing systems and practices. Specifically: 
	\begin{itemize}
		\item \textbf{Limited Accessibility}: The current system and accompanying practices are incapable of adjusting to the patient's needs and makes it difficult for the patient to determine said needs, limiting accessibility for first time or inexperienced users. 
		\item \textbf{Lack of Integration}: The current patient care coordination systems of the hospital  do not integrate into the respective clinician's scheduling and coordination systems, adding a barrier to the patient's appointment requests and confirmations.
		\item \textbf{Disjoint Systems}: The current means of clinician outreach for patients does not have a central, unified system.   
	\end{itemize}
	Alleviating these process hurdles through the design and development of a new care coordination system utilizing HCD may enhance current clinical workflows and improve patient outcomes. As such, this study aims to address the following research questions: 
	\begin{enumerate}
		\item What are patient experiences indicative of an ineffective experience within MMC's care coordination system? 
		\item How can HCD be leveraged to design and develop a care coordination system for MMC?
		\item How does stakeholder-driven development influence the success of a new care coordination system at MMC? 
	\end{enumerate}
	
	\subsection{Objectives of the Study}
	This study aims to design and develop a system to enhance the current patient care coordination processes of the MMC. To achieve this goal, the following are defined: 
	\subsubsection{General Objective}
	\begin{itemize} 
		\item To design and develop a care coordination system for the MMC leveraging HCD. 
	\end{itemize}
	\subsubsection{System Objectives}
	\begin{enumerate}
		\item \textbf{Clinician and Secretary}:
		\begin{itemize}
			\item \textbf{Clinician Profile Manager}: Through the interface, the user manages essential clinician data visible to the patient. 
			\item \textbf{Appointment Management}: Through the interface, the user manages their respective clinician's appointments. 
		\end{itemize}
		\item \textbf{Patient}:
		\begin{itemize}
			\item \textbf{Clinician Directory}: Provides patients with a clear and organized list of available clinicians alongside their respective schedules. 
			\item \textbf{Account Creation}: Guides patients through the account making process, enabling the user to book appointments, view appointment history, and manage their appointments.
			\item \textbf{Appointment Booking}: The patient encodes the necessary details and the appointment request is forwarded to the clinician and secretary via the preferred means. 
			\item \textbf{Appointment via QR}  (Kiosk): The system generates a QR Code for the patient to continue the appointment booking process on their mobile device.
			\item \textbf{Appointment Management}: Displays patient appointment history and enables management options
		\end{itemize}
		\item \textbf{System Administrator}:
		\begin{itemize}
			\item \textbf{Administrator Interface}: Administrative interface for control over the system and system monitoring
		\end{itemize}
	\end{enumerate}
	\subsubsection{Human-Centered Design Objectives}
	\begin{itemize}
		\item To design and deliver project artifacts, ensuring HCD is applied through adherence to the most recent Centre for eHealth and Wellbeing Research Roadmap. 
		\item To conduct active feedback collection involving stakeholder groups to ensure their input is a factor in the design and development of the system.
		\item To evaluate the system throughout development through Usability Testing and the User Experience Questionnaire. 
		\item To conduct a final evaluation of the system in coordination with stakeholder groups in conjunction with the User Experience Questionnaire.
	\end{itemize}
	\subsection{Significance of the Project}
	 The study's significance lies in the following: 
	\begin{enumerate}
		\item \textbf{Institute}: The resulting system, alongside the findings of the study,  will serve as a meaningful approach to addressing the institute's process hurdles in patient care coordination and widen the accessibility of its services to patients by enabling transparent communication and efficient coordination between the patients and clinicians. 
		\item \textbf{Users}: For clinicians and their respective secretaries, the findings of the study may address  process hurdles in relation to patient outreach. For patients, the findings of the study may contribute towards improving their health outcomes through widened accessibility and ease of communication with the institute's clinicians. 
		\item \textbf{Future Researchers}: The findings of the study will contribute towards future efforts targeting possible improvements towards the patient-clinician relationship through approaches that leverage HCD. 
	\end{enumerate}
	\subsection{Scope and Limitations}
	The following are the scope and limitations of the study: 
	\begin{enumerate}
		\item The system is designed to streamline the institute's care coordination process, specifically the patient's interactions with the institute prior to consultation with their clinician. The relevant processes consist of the patient seeking out their desired clinician, booking an appointment, and appointment confirmation. As such, the system will not handle tasks beyond the point of patient referral or appointment confirmation, such as diagnosing, billing, or post-consultation record management.
		\item Stakeholders for the system development and implementations will consist of clinicians, their respective secretaries, and patients. 
		\item The study will acquire necessary approval from the UP Manila Research Ethics Board (UPM REB), specifically documents under Standard Operating Procedures II  (SOP II), before collecting any responses and feedback from stakeholders and participants.
		\item The system will be designed to be architecturally agnostic and capable of servicing the diverse departments of MMC. However, in order to maximize focus, integration quality, and success probability, the development processes will initially interface with the Rheumatology Department.  
		\item The system is designed to be partially functional for patients without an account. Account creation is required to enable appointment booking, appointment history viewing and continuous interaction with the system for future visits. 
		\item The system will only handle non-diagnostic patient information necessary for registration, referral, and appointment scheduling. Sensitive clinical data such as patient test results and health records will not be collected or stored. 
		\item The project's implementation and evaluation primarily involves three user roles, Patients, Secretaries, and the Clinician. System administrators are only included for system maintenance and account management, administrator access to patient-related data is strictly limited to non-sensitive information. 
		\item The system will operate independently of existing institute information systems, electronic health record infrastructures, and health service queuing systems. Integration with external databases, payment processors, or telemedicine platforms or systems is outside scope.
		\item  The institute agrees to the means of evaluation for the system; furthermore, the users agree to have their feedback recorded to further along the design and development processes of the system. 
	\end{enumerate}
	
	\subsection{Assumptions}
	The following are the assumptions of the study: 
	\begin{enumerate}
		\item The institute will provide access to existing clinician schedule data and other non-sensitive information relevant for testing and populating the system during development and evaluation. 
		\item Secretaries, clinicians, and patients are assumed to be available and willing to participate in user acceptance testing activities, providing consent and feedback for system evaluation.  
		% \item The existing devices are sufficient and capable of running the system. Additionally, any server or similar deployment infrastructure will be handled by the institute
		\item Patients interfacing with the on-site kiosk currently have possession of their mobile device capable of scanning QR codes and rendering the web-interface. 
	\end{enumerate}
	
	
	\section{Review of Related Literature}
	
	%REVISE THIS TO FOLLOW NEW OUTLINE AND TACKLE EVALUATION 
	This chapter will cover a review of literature tackling the primary concepts and foundations behind the development of this system. The review will tackle literature discussing the value of triage and care coordination in clinical institutions followed by essential definitions of digital health, and the documented issues with Health-Information Systems (HIS) ranging from international to local issues with HIS in the Philippines. Subsequently, literature tackling definitions and applications of HCD in healthcare. The chapter aims to discuss the role of technology in a healthcare setting and the value it provides to professionals by curating prior works and reviews within the domain and identifying any challenges and limitations faced by respective authors. Additionally, the impact of HCD  applied in prior works and the methodology behind their evaluation are to be explored. The review concludes with a synthesis recounting the process burdens by the MMC and how HCD will be leveraged to address these burdens. 
	% REWRITE INTO PAR TALKING ABOUT CARE COORDINATION 
	\subsection{Triage and Care Coordination} 
	Triage is defined as the vital component in healthcare that encompasses the process of assessing and prioritizing patients with the aim of assigning acuity level and how long the patient can safely wait \cite{peta_2023_triage}. Although variations exist, it is important that the patient's first point of contact whether system or human has the capacity to place them in the right care unit as accurately as possible. Select reviews delve further into the factors that affect proper triage, stating that attitude, capabilities, and experience of personnel are vital, placing  strong emphasis on instilling awareness among staff and healthcare professionals ensuring a safe environment for the patients and increasing the chances of improved health outcomes \cite{fekonja_2023_factors, gorick_2022_factors}. 
	
	In the Philippines patients do not observe the practice of first consulting a primary care provider instead bypassing the primary care level and directly consulting with their specialist of choice. The lack of efforts towards providing a ``gatekeeper" within the patients' healthcare journey leads to significant inefficiencies within the process; consequently, the country faces issues with patients ``physician hopping". Such patient behavior is also attributed to the societal perception that family medicine is not a ``high-profile profession" and is associated with less expertise, pressuring institutes in urban areas towards hosting specialists for professional and financial viability \cite{Lajeunesse2012}. In light of this, the Philippine Department of Health \cite{DOH_AO_2020_0024} issued Administrative Order No. 2020-0024 which directly addresses the resulting systemic inefficiencies by establishing that the primary contact level must act as the navigator, coordinator, and initial and continuing point of contact within a patient's healthcare journey.
	
	As a result of the ``specialist-mindset" Filipinos fall into, patient-clinician communication falls up to the preferred means of communication of both parties such as Rakuten Viber or Text SMS. In a personal interview with Dr. Andrei Rodriguez \cite{Rodriguez_A}, they describe their system as one that has their secretary being the patient's first point of contact via the aforementioned means in regards to appointments and consultations, similarly for other Rheumatologists within the MMC . Such systems, while modern and up to societal norm, may lead to miscommunication, mishandling of appointments, and other hurdles stemming from the distributed human burden of said system. Chang et al. \cite{Chang682} discusses the diverse patient perspectives and how disjoint systems can be overwhelming for the patient and consequently lead to poorer health outcomes. Further discussing that patients often reported personally handling their own care coordination due to a lack of external support, with family members taking on the responsibility of scheduling visits and searching for providers .
	
	\subsection{Digital Health}
	Since 2020, Digital Health has been an emerging research domain, continuously ushering technological innovations and intersecting expertise between healthcare and technology. However, due to its broadness, a formal definition has been under debate among researchers in the field. A thorough review of published definitions of Digital Health revealed a majority of authors formally recognize Digital Health as a domain more concerned with health  over technology \cite{fatehi_2020_what}. With the latter subdomain concerned more on the proper usage and utilization of technologies to improve health outcomes. Essentially, Digital Health concerns itself with the intelligent use of technology to improve overall population well-being.
	
	\subsection{Issues with Health Information Systems}
	A Health Information System (HIS) is defined as critical systems deployed to help organizations modernize healthcare processes through intelligent integration. However, the current state of HIS deployments has been deemed inadequate due to circumstantial challenges that arise throughout numerous widespread efforts. Epizitone et al. \cite{epizitone_2023_a} held a review of literature documenting various HIS deployments in the health sector emphasizing the rich history of such systems and the various challenges that arise and how they were overcome. The evolution of these systems over time concerned itself with the transition from manual means to the integration of technology and the digitization of data and operations. Factors such as technological (clinical documentation, decision support), organizational (budget, teaching status), environmental (usage, training), and human (computer literacy, user friendliness) were documented as valuable input in the success or failure of the adoption of new systems in the sector. The authors then describe the value a well-implemented HIS presents to its stakeholders, these come in the form of ease of access, cost reduction, open communication, and the introduction of novel modes of healthcare delivery. 
	
	As technology continues to evolve, patients find themselves interfacing with more Digital Health Interventions (DHI) as time goes on. There exist DHIs catered towards enabling a support system for patients managing chronic illnesses, such products alongside other public-facing DHIs are seen to be viable means to promote preventative healthcare. A second review of qualitative studies found that, alongside the evolution of the technology itself, barriers arise preventing people from interfacing and evaluating public-facing DHIs. Primarily, these barriers come in the form of personal agenda and motivation, personal life and values, engagement and recruitment approaches, and lastly the perceived quality of the DHI in question \cite{oconnor_2016_understanding}. The findings emphasized the user’s perceived value and the amount of trust they can place in the system as a consequence are vital in ensuring their participation and evaluation of DHIs. 
	
	In developing nations, the deployment of a sustainable HIS is consistently hindered by outdated technology, poor infrastructures, and insufficient government policies. Rahman et al. found that countries like Bangladesh find its patients receiving ineffective healthcare due to technological factors such as fragmented and non-digitized data. Compounding the difficulties faced are barriers in implementations such as education, infrastructure, financial support, and cultural and political issues respective to the nation. The authors then explored these barriers in detail, finding that the shortcomings in education lead to diminished enthusiasm, minimized adoption rates, and poor population awareness for the system. Shortcomings in infrastructure manifest in inaccessible networks and absence of centralized healthcare databases; while cultural and political issues manifest in high unfamiliarity and low acceptance rates for the technologies \cite{rahman_2024_healthcare}. 
	
	Throughout Southeast Asia, the primary challenge of HIS efforts is to reach the unreached population regarding the accessibility of systems. Lu et al. \cite{lu_2021_assessment} discusses the key challenges for systems in these nations, which primarily include lack of policy and strategy, uncoordinated investment in respective technologies, low cooperation and capabilities, and poor overall infrastructure ranging from communication to interoperability.  In local context, the Philippines faces the very same barriers in HIS efforts, alongside a lack of policies and additional challenges for system development efforts. In essence, the Philippines has yet to fulfill its potential to deploy a complete HIS. Garcia et al. \cite{garcia_2021_health} found that instances of underutilized digital patient records, and over-complicated data extraction are what hinders the Philippines from deploying systems that can cater not only to the elderly, but the overall population. 
	
	The recurring challenges discussed illustrate that there exists a considerable portion of the shortcomings of HIS that stem not only from the technological aspect, but from a notable lack of attention to human and contextual factors influencing use cases. As such, addressing the issues requires an approach that prioritizes stakeholder needs and accommodates a tailored user experience throughout the design and development of such a system. Leary et al. \cite{leary_2022_an} defines HCD as a rigorous, innovative approach to solving problems and creating systems targeted towards the direct needs of the stakeholders. The origins of HCD date back to the 1960s, where a synthesis finds that the creation of DHIs that primarily incorporate the needs of the people is a vital step towards improving health outcomes. Additional findings reveal that DHI providers do not universally apply the same approach, and that not all phases of HCD are used equally. However, the core of each approach always keeps the stakeholders in mind throughout every step of the process. Human-Centered Design builds the theoretical foundation for the remainder of the chapter, further exploring how the approach has been applied in healthcare settings to support adoption, usability, and patient engagement. 
	
	
	\subsection{Human Centered Design in the Health Sector}
	Göttgens et al. \cite{gttgens_2021_the} explores the current practices of HCD towards the development of innovations in the field, and throughout each implementation they found that while HCD fails to conclude with a single correct answer, they enable active engagement with stakeholders. For instance, Watchler et al. \cite{wachtler_2018_development} aimed to address patient engagement in an app targeted towards patients of depression. Utilizing an iterative user-centered design (UCD) process. Involving two focus groups where user input was processed and used by their computer engineering team to develop prototypes of the application. The authors found that the UCD approach led to valuable insight and improvements towards the development of their application, stating that users felt hope and engaged in self-reflection. Woodard et al. \cite{woodard_2018_the} sought to improve cancer survivors' awareness of fertility preservation by developing a decision support system for these women. For development, they utilized a UCD involving groups of clinicians and patients. Afterwards post-decision scores amongst survivors were assessed, determining the system's feasibility. The authors found the iterative cycles of user evaluation aided the development team in meeting patients and survivors needs inclusive of patient friendly information and inclusion of novel features. Compounded by the high scores on the implemented knowledge preservation scale.
	
	Eberhart et al. \cite{eberhart_2019_using} explored the utility of HCD in improving health outcomes for patients with asthma, and introduce HCD as a multiphase process aimed at identifying the needs of end users and develop solutions around them. Their utilization of HCD consisted of ethnographic-style observations which served as a basis for user experience and was used to determine achievable areas regarding application development. Leading to the development of a successful prototype that involved a diverse population of stakeholders, and leaves opportunity for more community involved iterations. Catalani et al. \cite{catalani_2014_a} aimed to integrate patient care for those diagnosed with HIV and Tuberculosis, employing an approach that concerned itself more on stakeholder interviews and laboratory simulations and in-context usability testing. Their approach allowed them to successfully implement a well-informed system in a self-described complex and resource constrained context. Human-Centered Design sees its utility in complex clinical contexts, evidenced by Wentzel's \cite{wentzel_keeping_an_eye} work on infection control and antimicrobial stewardship. The author cited that focusing on the direct needs of stakeholders ensures that the technology design matches user abilities and preferences, promoting better usability and effectiveness. Efforts that take stakeholders into account are necessary in order to overcome the shortcomings of HIS deployments that neglected human and contextual factors. 
	
	
	Formalizing a developmental approach is necessary to ensure system viability and successful deployments, to address this, studies have anchored their methodologies on the Centre for eHealth and Wellbeing Research (CeHRes) Roadmap \cite{Kip2025CeHRes} to guide their development and evaluation processes to ensure they meet the needs of stakeholders. Wentzel's \cite{wentzel_keeping_an_eye} work on the Antibiotic App and MRSA-net was informed by the CeHRes roadmap and employed methods such as focus groups and card sorting to discover user mental models, translate ``top-down" medical protocols into ``bottom-up" task-oriented support, and identify the burdens associated with scattered information infrastructure. The holistic nature of the roadmap enabled iterative development phases that prompted stakeholder discussions which further informed the development of the system. The CeHRes roadmap enhances HCD by providing a holistic framework that enables integration of HCD activities with business modeling throughout the development period, ensuring that the stakeholders are serviced and the system is suitable given organizational contexts \cite{wentzel_keeping_an_eye,Kip2025CeHRes,vanVelsen2012Personas}.
	
	% PARA ON EVALUATION 
	Consequently, the evaluation of a Health Information System is a crucial step to ensure the project meets as much of the stakeholder needs. Additionally, evaluation serves as a means to manage unwanted outcomes such as functional errors, lack of user-friendliness, and system unreliability. Kia et al.\cite{sharifikia_2022_health} emphasizes the importance of continuous evaluation of modern health technologies, further discussing that despite the volume of literature on HIS, there is a need for a standardized tool for the assessment and evaluation of these technologies. Recognizing this, Kushendriawan et al.  \cite{aldikushendriawan_2021_evaluating} employed the User Experience Questionnaire as an easy-to-apply, valid, and reliable means to acquire a subjective quality assessment of a product. The User Experience Questionnaire is an evaluation means introduced by Laugwitz et al. \cite{laugwitz_2008_construction} as a means to measure user experience by considering task-related usability goals and non-task-related user experience goals dependent on six scales: Attractiveness, Perspicuity, Efficiency, Dependability, Stimulation, and Novelty.

	The study operates within the domain of care coordination and digital health within the MMC, a leading Philippine institution. The burden of process hurdles in the MMC's patient care coordination stems from disjointed systems suffering from ineffective user experience, limited accessibility, and a lack of integration. Such local issues echo broader  challenges in the adoption of Health Information Systems concerned with human, organizational, and technological factors. The study will leverage HCD to address this as validated by the discussed literature, citing enhanced patient engagement, improved usability, and the affordances provided by a robust, tailored digital health intervention that is well-informed on contextual factors and stakeholder needs. 

	
	
	\section{Theoretical Framework}
	
	\subsection{Makati Medical Center}
	The Makati Medical Center is among the top healthcare providers in the Philippines accredited by the Department of Health, United Nations Children's Fund, World Health Organization, Joint Commission International among others \cite{makatimedicalcenter_about, makatimedicalcenter_2025_makatimed}. Hosting a wide gamut of well-trained and board certified staff dedicated towards improving health outcomes \cite{kritikakritika_2025_top, metroguide_2021_the, flor_2023_top}. 
	
	The MMC is categorized as a Tertiary Private Training Hospital with a 600-bed operating capacity \cite{MMC_Proceedings_Mar2025_Vol20, MMC_Proceedings_Sep2023}. Hosting a wide gamut of facilities such as a 50-bed Emergency Department complex alongside other critical care units such as a Pediatric Intensive Care Unit and a Neonatal Intensive Care Unit among others \cite{MMC_Compendium_2022_2023}. The institution maintains its commitment to education through specialty society-accredited internship, residency, and sub-specialty fellowship training programs, designed to train succeeding generations of healthcare professionals \cite{MMC_Compendium_2022_2023}. 
	
	%\subsection{Web Platforms}
	\subsection{Human Centered Design}
	Human-Centered Design (HCD) serves as a core theoretical framework for developing rigorous, innovative, and ``human-centered" solutions \cite{leary_2022_an, gttgens_2021_the}. Within the context of this study, the approach will be leveraged to address the local issues at MMC reflective of the broader adoption challenged faced by the deployment of Health Information Systems (HIS), specifically issues concerned with human, organizational, and technological factors. Through the approach, this study aims to design and develop a robust, tailored digital health intervention streamlining the patient's interactions with the institute prior to clinician consult. 
	
	The application of HCD is primarily characterized by collaborative, creative, and iterative mindsets rooted in empathy with respective stakeholders in order to ensure that the developed solution aligns with user mental modes and existing work practices \cite{leary_2022_an, gttgens_2021_the, wentzel_keeping_an_eye}. Methodologies associated with HCD are essential for overcoming potential mismatches between technology features and the actual needs of the user, typically involving ongoing and continuous evaluation \cite{wentzel_keeping_an_eye}. Ultimately, the active involvement of stakeholders throughout the design and development processes positions the study to create a successful digital health intervention that addresses the challenges discussed in accordance with HCD.
	
	
	\subsection{The CeHRes Roadmap 2.0}
	The Centre for eHealth and Wellbeing Research (CeHRes) Roadmap \cite{Kip2025CeHRes} is a holistic and interdisciplinary framework designed to guide the development, implementation, and evaluation of eHealth systems in order to ensure that the resulting system meets the needs of the intended users and stakeholders within the context it will be used in. The roadmap structures the development processes into five (5) highly interrelated phases: \textbf{Contextual Inquiry}, \textbf{Value Specification}, \textbf{Design}, \textbf{Operationalization}, and \textbf{Summative Evaluation}, each connected by continuous evaluation cycles.
	
	The updated and most recent version of the CeHRes roadmap will serve as the guiding framework for this study. Figure \ref{fig:cehres_rm}, taken from the publication of Kip et al. \cite{Kip2025CeHRes}, details the roadmap alongside the five guiding pillars for the approach . Specifically: 
	\begin{enumerate}
		\item \textbf{Ongoing and Intertwined}: The development, implementation, and evaluation of an eHealth system are ``ongoing and intertwined" processes that are interrelated and relevant from the start. 
		\item \textbf{Holistic}: The development, implementation, and evaluation of an eHealth system require a holistic approach, concerned with the relationships present between technology, people, and context. 
		\item \textbf{Continuous Evaluation}: The development, implementation, and evaluation of an eHealth system benefits from an iterative approach, characterized by continuous evaluation cycles ensuring coherence between the outcomes of each phase of development and alignment with stakeholder needs 
		\item \textbf{Stakeholder Involvement}: To successfully develop an eHealth system that meets the needs and accommodates all stakeholders, their active and consistent involvement is essential throughout the entire process.
		\item \textbf{Interdisciplinary}: The development, implementation, and evaluation processes of an eHealth system are based on interdisciplinary approaches that combine and integrate methods and models from differing disciplines. 
	\end{enumerate}
	% Roadmap 
	\begin{figure} [htbp]
		\centering
		\includegraphics[width=0.9\linewidth]{draft-figures/CeHRes Roadmap.jpg}
		\caption{Updated CeHRes Roadmap, taken from Kip et al.\cite{Kip2025CeHRes}}
		\label{fig:cehres_rm}
	\end{figure}
	
	\subsubsection{Contextual Inquiry}
	The initial phase laid out in the CeHRes roadmap, Contextual Inquiry serves as the means to document the environment, the stakeholders, and the current situation with the aim of identifying key areas for improvement and where the system offers its value. The phase focuses attention on the stakeholders and their context from the onset of the development process, increasing the chances of a specific, context-relevant solution. The outcome of which provides the foundation for all succeeding development, implementation, and evaluation activities \cite{Kip2025CeHRes}.
	
	\subsubsection{Value Specification}
	The Value Specification phase involves systematically narrowing down the topics identified in the Contextual Inquiry phase and translating them into tangible requirements that build the foundation for the content and design of the resulting system. The phase entails formulating and prioritizing values from key stakeholders and using this information to set concrete goals and specific requirements descriptive of the technology \cite{Kip2025CeHRes}. 

	\subsubsection{Design}
	The Design phase is the iterative, dynamic, creative, and collaborative period where the requirements established from the Value Specification phase are to be translated into visual representation through prototyping of the resulting system. The phase aims to facilitate active collaboration between stakeholders and the developer, enabling constant testing and updating of the prototypes, ensuring that the technology is refined before full implementation \cite{Kip2025CeHRes}.
	
	\subsubsection{Operationalization}
	The Operationalization phase focuses on the planning and execution necessary for the introduction, dissemination, adoption, and internalization of the first functional version of the system within its operational context. This phase encompasses the formal launch of the technology and organizational procedures are put into practice. A crucial step in this phase is to identify implementation barriers and facilitators which influences the eventual implementation plan \cite{Kip2025CeHRes}.
	
	\subsubsection{Summative Evaluation}
	The Summative Evaluation phase is an iterative, multi-method process highly intertwined with the development and implementation activities. The primary objective is to determine to what extent the system successfully meets the objectives of the previous phases and addresses the objectives of the study. The phase covers three main aspects: determining the impact of the technology within contextual applications, assessing the uptake of the system in terms of adoption and use, and finally investigating the relevant working mechanisms that explain whether or not the system was effective \cite{Kip2025CeHRes}. Ultimately, the evaluation provides insight into identifying further points of improvement in the implementation strategy. 
	
	\subsection{Usability Testing}
	Usability Testing (UT) is a fundamental method in the development of systems grounded on HCD. Defined as a process where performance and subjective data on a product is collected through representative users performing realistic tasks within the system \cite{lewis2012usability}. Within the context of this study, UT will serve as a means to conduct essential, iterative evaluation activities. Ultimately, UT ensures that the system is easy to use and acceptable among the stakeholder groups; furthermore, UT will guide the development of the system towards making improvements to maximize system effectiveness, efficiency, and overall user satisfaction.  
	\subsection{User Experience Questionnaire}
	The User Experience Questionnaire (UEQ) was developed as an end-user questionnaire aimed at measuring user experience in an efficient, quick, and simple manner; allowing users to express feelings, impressions, and attitudes immediately after experiencing a product. The questionnaire itself is structured as seven-stage semantic differentials, where each of the 26 items utilizes bipolar terms that span a semantic dimension \cite{laugwitz_2008_construction}. The UEQ is based on a theoretical framework aimed at distinguishing between perceived pragmatic quality and perceived hedonic quality, where each item loads on to six distinct scales: 
	\begin{enumerate}
		\item \textbf{Attractiveness}: The user's general impression of the product, reflecting whether they like the product or not \cite{laugwitz_2008_construction}.
		\item \textbf{Efficiency}: The user's impression that their goals are reached quickly and efficiently via the system, indicative of clarity in interface organization \cite{laugwitz_2008_construction}. 
		\item \textbf{Perspicuity}: The user's impression on how easy the system is to learn to use \cite{laugwitz_2008_construction}.
		\item \textbf{Dependability}: The user's feelings about the safety and controllability of their interaction with the product \cite{laugwitz_2008_construction}. 
		\item \textbf{Stimulation}: The user's impression that it is interesting and fun to use the product \cite{laugwitz_2008_construction}. 
		\item \textbf{Novelty}: The user's impression that the product design is innovative, creative, and successfully captures their attention \cite{laugwitz_2008_construction}. 
	\end{enumerate}
	
	The UEQ is a well-established practice in software development cycles, as such the complete UEQ toolkit \cite{hinderks_user} will be utilized within the context of this study to streamline the evaluation process and ensure it adheres to an established standard. The following sections will detail the processes involved in the toolkit and their contributions.
	
	\subsubsection{Data Collection}
	The UEQ utilizes a \textbf{seven-stage differential} scale across all \textbf{26} of its items. Refer to Table \ref{tab:ueq_scales} for the item distribution, Figure \ref{fig:ueq_scales} for the scale structure of the questionnaire, and Figure \ref{fig:ueq_items} for the questionnaire itself. 
	\begin{table}[htbp]
		\centering
		\begin{tabular}{|l|p{6cm}|c|}
			\hline
			\multicolumn{1}{|c|}{UEQ Scale} & \multicolumn{1}{c|}{Measurement Focus} & No. of Items \\ \hline
			\textbf{Attractiveness} & General attitude towards the system & \textbf{6} \\ \hline
			\textbf{Perspicuity} & Ease of learning and understanding & \textbf{4} \\ \hline
			\textbf{Efficiency} & Impression that goals are accomplished quickly and efficiently & \textbf{4} \\ \hline
			\textbf{Dependability} & Feelings about system safety and controllability & \textbf{4} \\ \hline
			\textbf{Stimulation} & Impression that system is interesting and fun to use & \textbf{4} \\ \hline
			\textbf{Novelty} & Impression that system design is innovative and creative & \textbf{4} \\ \hline
		\end{tabular}
		\caption{Six Scales of the UEQ \cite{laugwitz_2008_construction,hinderks_user,schrepp_2023_user}}
		\label{tab:ueq_scales}
	\end{table}
	
	\begin{figure} [htbp]
		\centering
		\includegraphics[width=0.9\linewidth]{draft-figures/scalestruct.png}
		\caption{Scale Structure of the UEQ, taken from the UEQ Handbook \cite{schrepp_2023_user}}
		\label{fig:ueq_scales}
	\end{figure}
	
	\begin{figure} [htbp]
		\centering
		\includegraphics[width=0.9\linewidth]{draft-figures/UEQ ITSELF.png}
		\caption{User Experience Questionnaire, taken from Hinderks et al. \cite{hinderks_user}}
		\label{fig:ueq_items}
	\end{figure}
	\clearpage
	\subsubsection{Data Processing}
	The analysis of raw collected data will be done via the UEQ Excel Data Analysis Tool \cite{schrepp_2023_user}. The transformation process consists of converting the raw ratings (ranging from $+3$ to $-3$ on the scale) into structured measures for interpretation.
	
	Subsequently, the UEQ Excel Data Analysis Tool \cite{schrepp_2023_user} implements the Inconsistent and Identical Answers Heuristics to handle suspicious responses indicated by the distance between the best and worst answers exceeding 3 for 2 or 3 of the scales, and $>$15 identical answers respectively. 
	\subsubsection{Calculations}
	The UEQ Excel Data Analysis Tool \cite{schrepp_2023_user} automatically calculates the following metrics for the six UEQ Scales: 
	\begin{itemize}
		\item \textbf{Scale Means}: The mean value ($\bar{x}$) for each scale is calculated, with values greater than $+0.8$ indicative of a \textbf{positive evaluation}. Conversely values below $-0.8$ are indicative of a \textbf{negative evaluation}, and values in between constitute a \textbf{neutral evaluation}.
		%\[
		%\bar{x}=\frac{\sum x}{n}
		%\]
		%Where $\sum x$ is the sum of individual items $x$ within a scale, and $n$ is the total amount of items within the scale. 
		\item \textbf{Level of Agreement}: The standard deviation ($\sigma$) for each scale is calculated and the following interpretations are made \cite{schrepp_2023_user}: 
		\begin{enumerate}
			\item \textbf{High Agreement}: $\sigma <0.83$
			\item \textbf{Medium Agreement}: $0.83 < \sigma <1.01$
			\item \textbf{Low Agreement}: $\sigma > 1.01$
		\end{enumerate}
		\begin{figure}[htbp]
			\centering
			\includegraphics[width=0.9\linewidth]{draft-figures/goodproduct.png}
			\caption{Sample plot of a system with ``good" results, taken from the UEQ Handbook \cite{schrepp_2023_user}, plotting }
			\label{fig:gooodproduct}
		\end{figure}
		\item \textbf{Cronbach's Alpha ($\alpha$)}: Primarily viewed within the Tool itself, this metric is calculated to measure the consistency of each scale, ensuring all items within the scale measure a similar construct. Values greater than 0.6 or 0.7 are sufficient within the context of the UEQ.  
		\[
		\alpha= \frac{nr}{1+(n-1)r}
		\]
		Where $r$ is the mean correlation of the items in a scale and $n$ is the number of items in a scale
		\item \textbf{Benchmarking}: The system's evaluation is further contextualized through comparison against the UEQ benchmark dataset \cite{hinderks_user}, further clarifying the system's user experience quality.
		\begin{figure}[htbp]
			\centering
			\includegraphics[width=0.7\linewidth]{draft-figures/benchmark.png}
			\caption{Sample plot of a system with ``good" results compared against the UEQ Benchmark dataset, taken from the UEQ Handbook \cite{schrepp_2023_user}. Scale Means are the metrics for comparison. }
			\label{fig:benchmark}
		\end{figure}
	\end{itemize}
	
	
	\section{Design and Implementation}
	% From the problem definition, you identify key techniques and approaches you feel are necessary to successfully find a resolution. You can further motivate the approaches by outlining how they were used in other contexts. This is an extension of the Literature review but with emphasis on how you intend to apply the methodology/ies to your problem. The ERDs, DFDs and technical architecture are placed here.
	\subsection{Implementation Strategy}
	In accordance with the CeHRes Roadmap \cite{Kip2025CeHRes}, this study will divide its strategy into 5 key phases, mainly: \textbf{Contextual Inquiry}, \textbf{Value Specification}, \textbf{Design}, \textbf{Operationalization}, and \textbf{Summative Evaluation}. Ensuring the project is well-positioned to deliver a system well-informed on the needs of its stakeholder groups. 
	\begin{figure}[htbp]
		\centering
		\includegraphics[width=0.9\linewidth]{draft-figures/CeHRes Artifacts Full.png}
		\caption{The CeHRes Roadmap 2.0 modified from Kip et al. \cite{Kip2025CeHRes} to contain the adjustments made with respect to contextual application and the target deliverables per phase of the study.}
	\end{figure}
	\subsubsection{Contextual Inquiry and Value Specification}
	Defining the user requirements is a vital first step in involving the stakeholders in the design and development processes of the overall system. Adhering to the CeHRes Roadmap, the first phase of the study entails gathering information around the environment and context of the proposed system, the goal of which is to collect broad foundational input to define the initial requirements for the first prototype, and progress towards refinement and eventual completion. As such the following actions are discussed: 
	\begin{enumerate}
		\item \textbf{User Requirements Definition}: In order to build a basis for tangible system requirements for the first prototype, a semi-structured survey via Google Forms comprised of detailed, open-ended questions that seek insight regarding instances of ineffective user experience, limited accessibility, among the other process hurdles introduced. In addition, feedback to questions on existing systems, frequency of visits, and preferred means of communication are collected to establish measurable data on user needs. Lastly, upon consent, one-on-one interviews with stakeholders is to be conducted to deepen understanding of the existing work practices and mentalities. 
		\item \textbf{Iterative Feedback Collection}: As means of capturing user feedback, the Simple Remote Exploratory Test (SRET) is conducted alongside every high-fidelity iteration of the system, gathering contextual feedback on the current iteration. 
	\end{enumerate}
	\clearpage
	\subsubsection{Design}
	%The CeHRes Roadmap describes the Design phase as the iterative development and testing of low-fidelity and high-fidelity prototypes \cite{Kip2025CeHRes}. 
	Within the context of this study, this phase will consist of the brunt of the development of the system. Adopting the Agile \cite{shore2021art} approach for the technical development of the system, this section will detail the actions done during this phase of the project \\
	\textbf{Core Development and Iterative Refinement}: Consisting of the primary development of the system, after collection of stakeholder input, responses are translated into tangible requirements iteratively until stakeholder needs are met. Consequently, after each developmental iteration of the system, immediate user feedback is collected and translated into further tangible requirements for the system via the SRET. This may consist of bugs or issues to address as time progresses. Furthermore, this study will operate under the following definitions for respective prototypes during this phase: 
	\begin{itemize}
		\item \textbf{Low-Fidelity Prototype}: The primary deliverables for these iterations will consist of static mock-ups of the system complemented by the figures representing key workflows such as Appointment Booking and Management. These prototypes will be evaluated via a Cognitive Walkthrough by representatives from the stakeholder group for validation, this would comprise the first quarter of development (See Figure \ref{fig:gantt} for project Gantt Chart)
		\item \textbf{High-Fidelity Prototype}: The primary deliverables for these iterations will be the system prototype as it continues to gain functionality where further refinement will be driven by a form of Unmoderated Remote Usability Testing, specifically SRET \cite{lewis2012usability, usertesting_2019_17}, this would comprise the remaining period of development (See Figure \ref{fig:gantt} for project Gantt Chart)
	\end{itemize}
	Ultimately for either prototype, feedback is captured and reviewed in order to translate it into actionable revisions for the current state of the system, furthermore the feedback will aid in determining priorities and documenting findings to progress system development. 
	\clearpage 
	\subsubsection{Operationalization}
	% final PROTOTYPE 
	%The CeHRes Roadmap describes the Operationalization phase as the planning for the introduction, dissemination, adoption, and internalization of a first functioning version of the system \cite{Kip2025CeHRes}. 
	Within the context of this study, this phase encompasses the last stages of development after accounting for the stakeholder input for all iterations, resulting in the final prototype of the system. As such, this section will discuss the following actions: 
	\begin{itemize}
		\item \textbf{Prototype Finalization}: Towards the last leg of the development period (See Figure \ref{fig:gantt} for project Gantt Chart), a final and stable prototype incorporating all validated requirements will be set as the system's Minimum Viable Product (MVP). As such, the following functionalities must be validated in addition to further requirements that arise during development: 
		\begin{itemize}
			\item Clinician Directory and Patient Account Creation 
			\item End-to-End Appointment Booking 
			\item Schedule and Appointment Management Interface (Secretary and Clinician)
		\end{itemize}
		\item \textbf{Beta Testing} (BT): Once MVP is set, it will be prepared for a controlled and isolated deployment simulating the targeted context ensuring a stable environment for evaluation without impact towards live hospital operations. 
		\begin{itemize}
			\item \textbf{Deployment Details}: The application stack will be deployed on a server with no connections to the hospital's main networks. Subsequently, relevant mock data will be utilized throughout this deployment to ensure no leakage of real sensitive data. 
			\item \textbf{User Training}: The training strategy will be a role-specific approach ensuring all parties are well-informed of the protocol before participation. Participants will be provided a self-service kit that goes over the task list and system guide, video demonstrations, system access link, and informed consent forms. Furthermore, participants will be reminded that their candid feedback on the system is encouraged and their data is handled accordingly. 
			\item \textbf{Testing Environment}: The testing environment for the BT will be deliberated in coordination with key stakeholder groups, ensuring consent from all parties and availability from the target department. 
			\item \textbf{Procedure and Documentation}: A document will be devised detailing the BT's scope, test cases, the mock data used, and the workflow for each user role. During this period, any bugs or requested changes are logged and iteratively fed back into the development cycle established in the previous sections. 
		\end{itemize}
	\end{itemize}
	\subsubsection{Evaluation}
	%The CeHRes Roadmap describes the Summative Evaluation phase as the process of investigating the systems impact, uptake, and working mechanisms to identify points of improvement \cite{Kip2025CeHRes}. 
	Within the context of this study, evaluation activities will occur throughout development up until the final stages of the system, adhering to the interdisciplinary nature of the roadmap and HCD as a whole. As such, the following means of evaluation are discussed: 
	\begin{itemize}
		\item \textbf{Simple Remote Exploratory Test} (SRET) \\
		This is an unmoderated, remote method of evaluation conducted throughout the development period (See Figure \ref{fig:gantt} for project Gantt Chart) aimed at allowing users to openly share their input on the current iteration via minimal tasks or high-level goals, providing quick feedback for system refinement. The SRET will primarily consist of the testing group being sent a development build of the system, where they will conduct Exploratory Testing. Exploratory Testing will be done via high level, open-ended goals (e.g. find clinician and book appointment, update clinician schedule, manage appointments, etc.), allowing users to follow natural mental models in order to capture meaningful qualitative data. The users will record their feedback via context-appropriate means and submitted through Google Forms.  
		\item \textbf{User Experience Questionnaire} (UEQ) \\
		The summative means for system evaluation, consisting of the end-user questionnaire implemented via Google Forms designed to measure users' experience (Refer to Figure \ref{fig:ueq_items}). The administration and subsequent evaluation of the system will follow the pipeline as seen in Figure \ref{fig:ueq_pipeline}, additionally participants will be encouraged to answer all items candidly, ensuring their original impression is conveyed in their response. 
		\begin{figure} [htbp]
		\centering
		\includegraphics[width=1\linewidth]{draft-figures/finalfinalpipeline.png}
		\caption{Evaluation Pipeline}
		\label{fig:ueq_pipeline}
		\end{figure}	
	\end{itemize}

	Collected responses to the UEQ will be handled appropriately and processed via the UEQ Toolkit \cite{hinderks_user}. Following this, findings are interpreted and documented in preparation for the final stages of the project. 
	
	Furthermore, as a means of stakeholder communication, the documented results and the final state of the system are to be presented to key stakeholder groups and their thorough, open-ended qualitative feedback on the project and its development will be documented as part of the findings of this study. 


	\subsection{Workflow Integration}
	In order to fulfill the system's primary goal of addressing the resulting process hurdles, integration into the existing workflows among stakeholder groups is crucial in ensuring that no unnecessary burden is added. Through coordination and discussion with domain experts \cite{Rodriguez_A, Rodriguez_AG}, this section outlines the workflow-adjacent components of the system. 
	\subsubsection{Appointment Management Protocols}
	As a means to enhance the current appointment workflow within the MMC, the appointment component will serve as one of the core offerings of the system, providing a systematic means for bridging patients and clinicians of the MMC.
	\begin{itemize}
		\item \textbf{Clinician Profile Management} \\
		In order to address hurdles regarding current disjoint systems and communications, the secretary will assume the responsibility of centralizing the clinician's profile management inclusive of schedule updates and profile details. The added responsibility is a vital component in achieving overall process coherence and patient awareness. Through proactive schedule management, the system ensures patients are aware of clinician availability.\\
		\item \textbf{System Accommodation of Schedule Changes} \\
		The system will be designed in such a way that the volatile nature of stakeholders' schedules may be accommodated. In instances where the patient may need to have the appointment rescheduled, they may request to do so via the system. Similar accommodations are made for the clinicians themselves, should the clinicians schedule change. In either scenarios both parties are notified at least 24 hours before the change to allow for further coordination. This is complemented by the inherent responsibility to communicate which future dates are feasible for the reschedule.
		\item \textbf{Clinician Protection Against ``No-Show" Appointments} \\
		To avoid situations where the clinician may face a ``no-show" appointment, the system will have a feature in place to generate and send an e-mail to the patient prompting them for confirmation. This is complemented by the inherent responsibility of the secretary to confirm the appointment via Text SMS. 
		\item \textbf{Clinician Protection Against Fraudulent Appointments} \\
		To avoid situations where the clinician may face fraudulent appointments or appointments made maliciously, the system inherently requires that a user be registered before they are able to book an appointment. Additionally, the system will have internal protections in place to prevent duplicate entries from being permitted.
	\end{itemize}
	
	\subsubsection{Triage Implementation}
	As a means to enhance the current triage workflow within the Makati Medical Center, this component of the system is representative of a strategic digitization effort aiming to augment the existing manual triage done currently via the concierge. This section outlines the structure and functionalities of the digitized triage system and its refinement process. 

	The digitized triage workflow will be designed to seamlessly integrate with and enhance current the current human-led triage process, specifically addressing the role of Triage as the navigator, coordinator, and initial point of contact within the patient's healthcare journey \cite{Lajeunesse2012}. As such, the digitized triage will apply a specific, logical flow ensuring patient routing and referral, focused on accessibility and minimizing complexity. The following key steps are outlined: 
	\begin{enumerate}
		\item \textbf{Patient Status Assessment}: The system initially queries the patient on whether it is their first visit or not, establishing whether the patient needs the full guided triage or the patient is aware on which clinician they are to inquire with. 
		\item  \textbf{HMO Accommodation}: Following the status check, the system queries the patient on their Health Maintenance Organization (HMO), providing necessary financial context for the referral process (See Table \ref{tab:hmo-accred} for MMC HMO Accommodations \cite{makatimedicalcenter_doctors})
		\item \textbf{Symptom Assessment and Referral}: The system queries the patient's current and non-identifiable concerns, the system then dynamically presents the list of referred clinicians, completing the triage and the patient may begin the appointment booking process (See Figure \ref{fig:activity-diagram} for Activity Diagram). 
	\end{enumerate}	
	\textbf{Stipulation}: The safe and successful implementation of the digitized triage is dependent on consistent corroborations with key stakeholder groups and domain experts enhanced by the iterative nature of HCD. As such, a core stipulation of the implementation is its collaborative development, involving domain experts to ensure the language is non-identifiable and sufficiently meaningful for initial assessment and routing. Furthermore, while general user-feedback regarding experience with the digitized triage is collected, the core routing logic is governed by the input and consensus of domain experts and key stakeholder groups. The general user feedback informs the user interface and user experience refinement, and safety and routing remain a non-negotiable metric validated by expert review.
	
	\begin{table}[htbp]
		\centering
		\begin{tabular}{|p{0.45\linewidth}|p{0.45\linewidth}|}
			\hline
			AMAPhil, Inc. & Health Plan Philippines Inc. (HPPI) \\ \hline
			Avega Managed Care & Insular Health Care (I-Care) \\ \hline
			BenLife Insurance Co. Inc. (formerly Star Health Care) & Asalus (formerly Intellicare) \\ \hline
			Carehealth Plus & Kaiser International Healthcare Group \\ \hline
			Carewell Health Systems, Inc. & Life \& Health HMP, Inc. \\ \hline
			Caritas Health Shield & Maxicare \\ \hline
			Cocolife Healthcare & Medicard \\ \hline
			Cooperative Health Management Federation & Medicare Plus, Inc. \\ \hline
			Dynamic Care Corporation & Medocare \\ \hline
			EastWest Healthcare & Optimum Medical \& Healthcare Services, Inc. \\ \hline
			eTIQA (formerly AsianLife) & Pacific Cross Health Care (formerly BlueCross Health Care) \\ \hline
			Flexicare & Philcare \\ \hline
			Fortune Life Insurance Co., Inc. & Value Care Health Systems, Inc. \\ \hline
			GetWell Healh Systems Inc. & Wellcare \\ \hline
			Health Maintenance Inc. (HMI) & \\ \hline
		\end{tabular}
		\caption{Current MMC HMO Accreditation \cite{Rodriguez_A,makatimedicalcenter_doctors}}
		\label{tab:hmo-accred}
	\end{table}
	\clearpage
	
	\subsection{Technical Architecture}
	The proposed system is multi-interface web-based application serving the target user groups. The frontend component of the system will be developed upon the React-Vite framework, ensuring fast performance and a responsive design across various devices. The backend component of the system will be developed using the Python programming language leveraging the Django framework to build the RESTful API which facilitates seamless and secure data exchange between client and the database built on PostgreSQL. The system will then be implemented onto Docker, to ensure that the main components are separated into their respective isolated environment. Adding a boundary of security to avoid cross-service compromises alongside establishing a main internal means of communication via container-to-container networking. 
	\subsubsection{Website}
	The Alagang MMC website will be the primary offering of the study, allowing the targeted users to interface with the system within their given context. As such the system is divided into four levels of responsibility (See Figure \ref{fig:ucd} for the System Use-case Diagram): 
	\begin{enumerate}
		\item \textbf{Secretary}: Their responsibilities are inherited into the system. Essentially, they are responsible for updating the schedule of their respective Clinician within the system. Additionally, they are charged with managing the received appointment requests for their respective clinicians. Any decisions made regarding patient appointment status are then returned to the patients.
		\item \textbf{Clinician}: The clinician as a user shares the same set of capabilities within the system as their respective secretaries; however, this is implemented from an administrative perspective. Essentially, any decisions or updates made by them are given system priority. 
		\item \textbf{Patient}: The patient responsibility within the system lies primarily within the web component, where they are responsible for booking and managing their appointments (See Figure \ref{fig:activity-diagram} for Activity Diagram).
		\item \textbf{System Administrator}: The user is charged with maintaining the system on a technical level, managing and auditing user profiles paired with system monitoring via audit trailing. 
	\end{enumerate}
	The following are the minimum hardware specifications in order to run the system's website component: 
	\begin{enumerate}
		\item \textbf{Operating System}:
		\begin{itemize}
			\item \textbf{Windows}: Windows 11 
			\item \textbf{MacOS}: macOS 10.15 Catalina or bewer
			\item \textbf{Linux}: Any distribution released after 2018
		\end{itemize}
		\item \textbf{Processor} (CPU): Intel Core i5 (10th Gen or bewer) or AMD Ryzen 5 (5000 Series or newer)
		\item \textbf{Memory} (RAM): 8GB Minimum, 16GB Recommended
		\item \textbf{Storage}: 256GB 
	\end{enumerate}
	The following are the minimum hardware specifications in order to run the system's mobile website component: 
	\begin{enumerate}
		\item \textbf{Operating System}:
		\begin{itemize}
			\item \textbf{Android}: Version 9.0 or later 
			\item \textbf{iOS}: Version 11.0 or later 
		\end{itemize}
		\item \textbf{Processor} (CPU): Quad-core mobile processor 
		\item \textbf{Memory} (RAM): 4 GB 
		\item \textbf{Internal Storage}: 16GB 
		\item \textbf{Web Browser}: Latest version of Google Chrome or Safari
		\item \textbf{Connectivity}: 4G LTE or reliable Wi-Fi 
		\item \textbf{Functionality}: Rear-facing camera capable of auto-focus
	\end{enumerate}
	\subsubsection{Kiosk}
	Figures \ref{fig:onsite_kiosk} and \ref{fig:onsite_kiosk_close} shows the intended kiosk the system will be implemented in. The on-site kiosk serves as the means for patients to interface with the system within hospital premises. Patient interaction with the device consists of interfacing with the system and upon referral and selection of desired clinician, the system will either: 
	\begin{enumerate}
		\item Direct the patient to the designated room in situations where the desired clinician is currently \textbf{present}.
		\item Generate a QR code for the patient to scan and continue the appointment booking process from their capable mobile device in situations where the desired clinician is currently \textbf{unavailable}.
	\end{enumerate}
	\begin{figure} [htbp]
		\centering
		\includegraphics[width=0.4\linewidth]{draft-figures/kiosk.jpeg}
		\caption{On-Site kiosk present at Makati Medical Center}
		\label{fig:onsite_kiosk}
	\end{figure}
	\begin{figure} [htbp]
	\centering
	\includegraphics[width=0.4\linewidth]{draft-figures/kiosk2.jpeg}
	\caption{Close up view of on-site kiosk present at Makati Medical Center}
	\label{fig:onsite_kiosk_close}
	\end{figure}
	\clearpage
	The following section details the hardware specifications of the kiosk: 
	\begin{enumerate}
		\item \textbf{Display}:
		\begin{enumerate}
			\item \textbf{Diagonal Size}: 24 inches 
			\item \textbf{Resolution} FHD (1920x1080)
			\item \textbf{Technology} Projected Capacitive (PCAP) Touch
		\end{enumerate}
		\item \textbf{Internal Computer}: 
		\begin{enumerate}
			\item \textbf{Operating System}: Windows 10 IoT Enterprise
			\item \textbf{Processor} (CPU):  Intel Core i5 i5-1235U
			\item \textbf{Memory} (RAM): 8GB 
			\item \textbf{Graphics}: Intel Iris Xe Graphics 
		\end{enumerate}
	\end{enumerate}
	
	\subsubsection{System Architecture}
	This section will detail the envisioned components of the proposed system as validated by respective domain experts \cite{Rodriguez_A, Nicolas_AJ} and are subject to change throughout the development period. Furthermore, the section will expound on previously discussed protocols and stipulations within the context of system implementations. 
	\\ \textbf{Patient}
	\begin{itemize}
		\item \textbf{Account Creation}: Before patients are able to book appointments with clinicians, they will need to register for an account within the system. Refer to Table \ref{tab:basic_info} for the basic information required for account generation. Their initial credentials are sent to them and the user is advised to update these as soon as possible. 
		\begin{table}[htbp]
			\centering
			\begin{tabular}{|l|l|}
				\hline
				Patient Name        & Active Email \\ \hline
				Active Phone Number & Birthdate    \\ \hline
				Gender              & Address      \\ \hline
			\end{tabular}
			\caption{Basic Information Required for Account Generation}
			\label{tab:basic_info}
		\end{table}
		\item \textbf{Appointment Booking}: Refer to Table \ref{tab:appt-info} for the basic information required for appointment booking. Before a patient proceeds to the clinician's respective appointment booking page, they are prompted with reminders regarding on-site appointments and preparations to make before proceeding to their appointment (e.g. come on time, prepare paperwork, etc.). Furthermore, to reinforce clinician protection the system will cache the patient's data the first time they make an appointment in order to prevent multiple bookings to the same clinician under the same name.
		\begin{table}[htbp]
			\centering
			\begin{tabular}{|l|l|}
				\hline
				Appointment Date & Time Slot \\ \hline
				Chief Complaint  & Booking Type\\ \hline
			\end{tabular}
			\caption{Basic Information for Appointment Booking}
			\label{tab:appt-info}
		\end{table}
		\item \textbf{Appointment Management}: Inclusive of appointment history, this interface provides the means for the patient to manage their appointments within the system. This may entail viewing the appointment status (accepted or rescheduled), canceling, or requesting a reschedule.
	\end{itemize}


	\textbf{Clinician and Secretary} 
	\begin{itemize}
		\item \textbf{Clinician Profile}: The secretary is charged with populating and updating the clinician's profile data within the system. Refer to Table \ref{tab:cp_data} for the data used to build the profile. Further interactions with this component are periodic updates to the schedule when necessary.
		\begin{table}[htbp]
			\centering
			\begin{tabular}{|l|l|}
				\hline
				Clinician Name    & Clinician Practice \\ \hline
				Clinician Room     & Clinician Schedule  \\ \hline
				HMO Accreditation & Misc. Information  \\ \hline
			\end{tabular}
			\caption{Data to build Clinician Profile. Note: ``Misc. Information" refers to clinician's history or background in addition to the display picture for the system.}
			\label{tab:cp_data}
		\end{table}
		\item \textbf{Appointment Management}: The appointment management interface will enable the users to manage the  respective appointments through accepting or requesting for reschedules. 

	\end{itemize}

	% DIAGRAMS
	\clearpage
	\subsubsection{System Diagrams}
	\begin{figure} [htbp]
		\centering
		\includegraphics[width=0.6\linewidth]{draft-figures/activity-diagram.png}
		\caption{Activity Diagram for the Appointment Booking Process within Alagang MMC}
		\label{fig:activity-diagram}
	\end{figure}
	\begin{figure} [htbp]
		\centering
		\includegraphics[width=0.9\linewidth]{draft-figures/realucds.png}
		\caption{Use Case Diagram for Alagang MMC}
		\label{fig:ucd}
	\end{figure}
	\begin{figure} [htbp]
		\centering
		\includegraphics[width=0.75\linewidth]{draft-figures/level-0-dfd.png}
		\caption{Level 0 Data Flow Diagram of Alagang MMC}
		\label{fig:l0-dfd}
	\end{figure}
	\clearpage
	\begin{figure} [htbp]
		\centering
		\includegraphics[width=0.9\linewidth]{draft-figures/level-1-dfd.png}
		\caption{Level 1 Data Flow Diagram of Alagang MMC}
		\label{fig:l1-dfd}
	\end{figure}
	\subsection{Ethical Considerations}
	The development and implementation of the proposed system will be governed by strict adherence to ethical principles, ensuring the protection of all stakeholders involved. The study is classified as presenting minimal risk to participants, primarily concerning potential risks to data privacy or the accidental disclosure of contact information. To mitigate these concerns, participants are explicitly instructed to avoid encoding real, identifiable clinical information during testing phases, and the system is designed to operate on an isolated network stack independent of the hospital’s primary infrastructure. In the event of any privacy breach, the matter will be immediately escalated to the relevant institutional officials. While the broader potential benefits of the study include streamlined clinical workflows and improved patient-clinician communication, participants serving as system testers receive a small token of appreciation as a direct, tangible benefit for their contribution to the research.
	
	Recruitment and the informed consent process are managed entirely by the researcher. Participants will be recruited via online platforms and social media, where they are directed to the Google Forms survey containing an informed consent section (see Figure \ref{fig:inf-cons}) at the very beginning. No individual can proceed with the survey or participate in the study without first reading and agreeing to these terms, which detail the nature of the research and the voluntary nature of their involvement. Toward the end of the survey, respondents may opt-in to further participation as system testers for the Beta Test or request more information via email. This process ensures that all involvement is based on full, explicit consent, and participants retain the right to withdraw from the study at any time without penalty.
	
	All collected data will be stored securely on the Google Forms platform for the full duration of the study and the subsequent thesis defense required for the fulfillment of the Computer Science program. Upon the completion of all academic requirements, the researcher will dispose of all electronic files and data through permanent deletion from the platform. 
	\begin{figure} [htbp]
		\centering
		\includegraphics[width=0.9\linewidth]{draft-figures/infcons.png}
		\caption{Informed Consent Section in Google Forms}
		\label{fig:inf-cons}
	\end{figure}
	
	\section{Expected Output/Timeline}
	% GANT TCHART 	
	\subsection{System Wireframes}
	This section will detail the initial mock-up of the system and is not representative of the final state of any of the discussed the prototypes 
		\begin{figure} [htbp]
			\centering
			\includegraphics[width=0.5\linewidth]{draft-figures/wf1_cp.png}
			\caption{Wireframe for Clinician Profile}
			\label{fig:wf1}		
		\end{figure}
		\begin{figure} [htbp]
			\centering
			\includegraphics[width=0.5\linewidth]{draft-figures/wf2_ma.png}
			\caption{Wireframe for Appointment Management. \textbf{Note}: Interface is inherently the same across both roles, with the difference being the options seen by the user}
			\label{fig:wf2}		
		\end{figure}
		\begin{figure} [htbp]
			\centering
			\includegraphics[width=0.5\linewidth]{draft-figures/wf3_ba.png}
			\caption{Wireframe for Appointment Booking}
			\label{fig:wf3}		
		\end{figure}
	
	\begin{landscape}
		\subsection{Project Timeline}
		% Screenshots of your final software and Gantt chart of your activities are placed here.
		Figure \ref{fig:gantt} shows the projected project timeline for Alagang MMC. 
		\begin{figure} [htbp]
			\centering
			\includegraphics[width=1.0\linewidth]{draft-figures/finalgantt.png}
			\caption{Gantt Chart for Alagang MMC Project Timeline}
			\label{fig:gantt}
		\end{figure}
	\end{landscape}
	
	\addcontentsline{toc}{section}{References}
	\bibliographystyle{ieeetr}
	\bibliography{biblio}
\end{document}
